\chapter*{Notation}
\addcontentsline{toc}{chapter}{Notation}
Symbols are used as below in every chapter unless the chapter says otherwise. Bold upper case is a matrix, bold lower case a vector; a symbol that has two meanings is listed with both, and the context says which is meant.
\begingroup\small\renewcommand{\arraystretch}{1.15}
\noindent\begin{longtable}{@{}>{\raggedright\arraybackslash}p{0.25\textwidth}>{\raggedright\arraybackslash}p{0.7\textwidth}@{}}
\toprule
Symbol & Meaning \\
\midrule
\endfirsthead
\toprule
Symbol & Meaning \\
\midrule
\endhead
\multicolumn{2}{@{}l}{\emph{Full-order model and dynamics}} \\
\midrule
$\bm K,\ \bm M,\ \bm C$ & stiffness, mass and damping matrices of the full-order model \\
$\bm u,\ \bm v,\ \bm a$ & displacement, velocity and acceleration vectors; $\bm f$ is the external load \\
$\omega,\ f,\ \zeta,\ \eta$ & circular frequency, frequency in hertz, damping ratio, loss factor \\
$\bm\Lambda$ & diagonal matrix of squared natural frequencies $\omega_i^2$ \\
$h,\ \Delta t$ & mesh size; time step \\
$E,\ \nu,\ \mu,\ \kappa$ & Young's modulus, Poisson's ratio, shear modulus, bulk modulus \\
$\beta,\ \gamma$ & parameters of Newmark's method (Chapter~\ref{ch:dyn}); $\beta$ is also the stiffness coefficient of Rayleigh damping $\bm C=\alpha\bm M+\beta\bm K$ \\
$\Omega$ & circular frequency of a harmonic excitation, $\bm f=\bm F_0e^{\mathrm i\Omega t}$ (Chapter~\ref{ch:dyn}) \\
\midrule
\multicolumn{2}{@{}l}{\emph{Elements and nonlinear mechanics}} \\
\midrule
$\bm f_{int}(\bm u)$, $\bm K_t$ & internal force and tangent stiffness of a nonlinear model \\
$\bm F,\ \bm E,\ \bm S$ & deformation gradient, Green--Lagrange strain and second Piola--Kirchhoff stress of the total-Lagrangian formulation (Chapter~\ref{ch:nonlin}); $\bm F_{ext}$ is the external load, and the context says which is meant \\
$\bm R(\bm u)$ & residual of the equilibrium equations, $\bm R=\bm F_{ext}-\bm f_{int}(\bm u)$ \\
$\bm D$ & elasticity matrix of the material (Chapters~\ref{ch:elem},~\ref{ch:nonlin},~\ref{ch:batch}) \\
$\bm B(\bm u_e)$ & strain--displacement matrix of an element (Chapters~\ref{ch:elem},~\ref{ch:batch}); in the contact chapters $\bm B$ is the gap map, and the context says which \\
$\partial\bm N/\partial\bm X,\ w_g,\ n_g$ & reference shape-function gradients, weight of Gauss point $g$ and number of Gauss points (Chapters~\ref{ch:elem},~\ref{ch:batch}) \\
\midrule
\multicolumn{2}{@{}l}{\emph{Contact}} \\
\midrule
$\bm B,\ \bm c$ & gap map and offset of the contact constraints, $\bm B\bm u\ge\bm c$ \\
$\bm\lambda,\ \bm U_\lambda$ & contact forces (Lagrange multipliers), $\bm\lambda\ge\bm 0$; basis for the contact forces \\
$\bm g,\ \bm g(\bm u)$ & gap (distance to the obstacle) in the contact chapters; $\bm g=\partial L/\partial\bm u$ in Appendix~\ref{ch:diff} \\
$c_p$ & positive scale of the complementarity function of the contact problem (Chapter~\ref{ch:contact}) \\
$\mathcal A_k$ & active set of contact constraints at iteration $k$ \\
\midrule
\multicolumn{2}{@{}l}{\emph{Reduction}} \\
\midrule
$\bm V,\ \bm U$ & reduction basis (columns span the reduced space); $\bm U$ is used for the primal basis in contact chapters \\
$\bm q,\ \bm u_r$ & reduced coordinates, $\bm u\approx\bm V\bm q$ \\
$\bm A_r,\ \bm M_r,\ \bm B_r$ & reduced matrices: $\bm V^T\bm K\bm V$, $\bm V^T\bm M\bm V$, reduced constraint matrix \\
$n,\ r,\ p$ & number of unknowns of the full model; size of the reduced model; primal basis size (also polynomial or perturbation order where stated) \\
$\sigma_j$ & singular values (snapshot matrix) or, in Chapter~\ref{ch:bt}, Hankel singular values \\
$\bm G,\ \bm w$ & ECSW training matrix and element weights \\
$\mathcal R,\ \mathcal E,\ \mathcal P$ & required, selected and passive sets of elements in ECSW (Chapter~\ref{ch:hyper}) \\
$\varepsilon,\ \mathrm{tol}$ & tolerance (of a truncation, a fit or a convergence test; the chapter says which) \\
$\bm P$ & mortar operator that gives the slave interface values from the master ones (Chapter~\ref{ch:cms}) \\
$n_e,\ N_a$ & number of elements; number of snapshot columns with active contact (Chapter~\ref{ch:hyper}) \\
\midrule
\multicolumn{2}{@{}l}{\emph{Input--output systems, frequency domain and bounds}} \\
\midrule
$\bm H(s),\ \bm H(\mathrm i\omega)$ & transfer function of a linear system and its frequency response (Chapters~\ref{ch:bt},~\ref{ch:freq},~\ref{ch:nonint}) \\
$\bm W_c,\ \bm W_o$ & controllability and observability Gramians (Chapter~\ref{ch:bt}) \\
$\bm\rho,\ \beta_{\mathrm{LB}}$ & residual of a reduced solution; lower bound of the smallest singular value of the system matrix used in the certified error bound (Chapter~\ref{ch:freq}) \\
\midrule
\multicolumn{2}{@{}l}{\emph{Differentiable layer (Appendix~\ref{ch:diff})}} \\
\midrule
$\bm f_\theta,\ \bm\theta$ & trainable additive correction and its parameters (Appendix~\ref{ch:diff}) \\
$\bm K_{\mathrm{eff}}$ & tangent of the corrected problem, $\bm K_T+\partial\bm f_\theta/\partial\bm u$ (Appendix~\ref{ch:diff}) \\
$\bm\lambda$ (Appendix~\ref{ch:diff}) & adjoint vector, $\bm K_{\mathrm{eff}}^T\bm\lambda=-\bm g$ (not the contact force of the same symbol) \\
$\bm N$ & Cholesky factor of a mass matrix, $\bm M=\bm N\bm N^T$ (Appendix~\ref{ch:diff}) \\

\bottomrule
\end{longtable}
\endgroup
